\DSource{0136}{47}
\hypertarget{AB01-PDF0136-P01}{}
\DLine{AB01-PDF0136-body-L001}{}{}{\Indent Dicit [auctor]: Nunc aliter haec explicabimus. Circa centrum E et diametrum AC, sit}
\DLine{AB01-PDF0136-body-L002}{}{}{circulus eclipticae ABC; circa centrum H, excentricus FMG; et diametrus AC per ambo}
\DLine{AB01-PDF0136-body-L003}{}{p. 71.}{centra transeat. Erit F punctum apogei in circulo parecliptico, G perigei. Initio locum Solis}
\DLine{AB01-PDF0136-body-L004}{}{}{in excentrico in M esse ponamus, et arcum FM excentrici, quem Sol iam percurrit, $30^\circ$ com-}
\DLine{AB01-PDF0136-body-L005}{5}{}{plecti; angulus FHM erit quoque $30^\circ$. Lineam EH, ambo centra coniungentem, iam vidimus}
\DLine{AB01-PDF0136-body-L006}{}{}{$2^p4'\,\qfrac{3}{4}$ esse. Semidiametrum HM excentrici et lineam EM ducamus, HM usque ad punctum}
\DLine{AB01-PDF0136-body-L007}{}{}{L quo cum perpendiculari LE convenit protrahentes. Triangulum HLE rectangulum est; eius}
\DLine{AB01-PDF0136-body-L008}{}{}{angulus LHE aequat angulum datum FHM; arcus super EL, ex circulo triangulum HLE}
\DLine{AB01-PDF0136-body-L009}{}{}{circumdanti, si circumferentia $360^\circ$ continet, $30^\circ$ complectetur, eiusque sinus $30^p$ quoque ex}
\DLine{AB01-PDF0136-body-L010}{10}{}{partibus quarum linea HE inter ambo centra 60 habet. Remanebit linea LH, quadrantem per-}
\DLine{AB01-PDF0136-body-L011}{}{}{ficiens, $51^p57'41''$, nam arcus complementaris LH $60^\circ$ continet. Sed ea ratione qua linea HE}
\DLine{AB01-PDF0136-body-L012}{}{}{inter ambo centra $2^p4'\,\qfrac{3}{4}$ complectitur, linea EL $1^p2'22''\,\qfrac{1}{2}$ (1) erit, et linea LH, quadran-}
\DLine{AB01-PDF0136-body-L013}{}{}{tem perficiens, $1^p48'2''$; itaque linea tota LM $61^p48'2''$. Triangulum MLE rectangulum est;}
\DLine{AB01-PDF0136-body-L014}{}{}{eius hypotenusam EM $61^p48'35''$ esse vidimus. Sed ratione qua lineae EM $60^p$ tribuuntur,}
\DLine{AB01-PDF0136-body-L015}{15}{}{EL est $1^p33'$ (2), et arcus super eam $0^\circ57'49''$, cum circulus triangulum HLE circumdans}
\DLine{AB01-PDF0136-body-L016}{}{}{$360^\circ$ sit. Arcus igitur AB eclipticae $29^\circ2'11''$ fere remanebit.}
\hypertarget{AB01-PDF0136-P02}{}
\DLine{AB01-PDF0136-body-L017}{}{}{\Indent Item Solem in puncto D excentrici esse ponamus, arcumque FD (3) $150^\circ$ complecti; reli-}
\DLine{AB01-PDF0136-body-L018}{}{}{quus arcus DG, inter locum Solis et perigeum, $30^\circ$ erit. Lineas EK et HD, quarum utraque}
\DLine{AB01-PDF0136-body-L019}{}{}{est sui circuli semidiametrus, ducamus; item perpendicularem ES. Triangulum HSE rectan-}
\DLine{AB01-PDF0136-body-L020}{20}{}{gulum est; latus EH, ambo centra coniungens, latus ES, angulus DHG (4) nota sunt; erunt}
\DLine{AB01-PDF0136-body-L021}{}{}{ergo nota latus HS et angulus HES (5). Cognoscuntur igitur quoque linea DG et hypotenusa}
\DLine{AB01-PDF0136-body-L022}{}{}{ED trianguli rectanguli ESD. Cum vero datus arcus DG et datus angulus GHD $30^\circ$ sint ut}
\DLine{AB01-PDF0136-body-L023}{}{}{diximus, eorumque sinus $30^p$ pariter, arcus quoque ES circuli triangulum rectangulum ESH}
\DLine{AB01-PDF0136-body-L024}{}{}{circumdantis complectetur $30^\circ$ ex partibus quarum totus ille circulus $360^\circ$ habet, et eius sinus,}
\DLine{AB01-PDF0136-body-L025}{25}{p. 72.}{nempe cathetus ES, $30^p$ quoque e partibus quarum linea EH, semidiametrus huius circuli, $60^p$}
\DLine{AB01-PDF0136-body-L026}{}{}{complectitur. Sed ratione qua lineae EH $2^p4'\,\qfrac{3}{4}$ tribuuntur, cathetus ES erit $1^p2'22''\,\qfrac{1}{2}$ (6);}
\par\noindent\begin{minipage}[t]{0.54\FullSourceWidth}\vspace{0pt}\centering\hypertarget{AB01-PDF0136-F01}{}\includegraphics[width=\linewidth]{AB01-PDF0136-F01.png}\end{minipage}\hfill\begin{minipage}[t]{0.43\FullSourceWidth}\vspace{0pt}\setlength{\GutterOffset}{0.5700000000000001\FullSourceWidth}
\DLine{AB01-PDF0136-body-L027}{}{}{qua re latus SH $1^p48'2''$ remanebit.}
\DLine{AB01-PDF0136-body-L028}{}{}{Semidiametrus HD excentrici $60^p$ conti-}
\DLine{AB01-PDF0136-body-L029}{}{}{net, de quibus si SH dematur, remanent}
\DLine{AB01-PDF0136-body-L030}{30}{}{lineae SD $58^p11'58''$; igitur hypote-}
\DLine{AB01-PDF0136-body-L031}{}{}{nusa ED trianguli rectanguli ESD erit}
\DLine{AB01-PDF0136-body-L032}{}{}{circiter $58^p12'32''$ (7). Sed ex partibus}
\DLine{AB01-PDF0136-body-L033}{}{}{quarum linea ED 60 habet, cathetus ES}
\DLine{AB01-PDF0136-body-L034}{}{}{$1^p4'17''$ complectetur, et arcus super}
\DLine{AB01-PDF0136-body-L035}{35}{}{eam, qui est quantitas anomaliae, $1^\circ1'$}
\DLine{AB01-PDF0136-body-L036}{}{}{$24''$ (8): arcus igitur KC eclipticae, $31^\circ$}
\DLine{AB01-PDF0136-body-L037}{}{}{$1'24''$ (9) circiter erit. Haec de anomalia}
\DLine{AB01-PDF0136-body-L038}{}{}{sufficiunt, quam explicare volebamus.}
\hypertarget{AB01-PDF0136-P03}{}
\DLine{AB01-PDF0136-body-L039}{}{}{\Indent Dicit [auctor]: Hunc motum ita ad}\end{minipage}\par
\DLine{AB01-PDF0136-body-L040}{40}{}{singulos gradus in tabulis a puncto apogei posuimus; similiter reperiuntur aequatio simplex}
\par\vspace{6pt}\hrule height.25pt\vspace{5pt}
\noindent\begin{minipage}[t]{.48\textwidth}\vspace{0pt}\fontsize{9}{11.5}\selectfont\setlength{\GutterOffset}{0pt}
\hypertarget{AB01-PDF0136-N01}{}
\DLine{AB01-PDF0136-notes_left-L001}{}{}{\Indent (1) Secundae apud Platonem $55''$.}
\hypertarget{AB01-PDF0136-N02}{}
\DLine{AB01-PDF0136-notes_left-L002}{}{}{\Indent (2) Cod. « unius partis et trium ac triginta se-}
\DLine{AB01-PDF0136-notes_left-L003}{}{}{« cundarum ».}
\hypertarget{AB01-PDF0136-N03}{}
\DLine{AB01-PDF0136-notes_left-L004}{}{}{\Indent (3) Cod. AD. Postea Plato 120.}
\hypertarget{AB01-PDF0136-N04}{}
\DLine{AB01-PDF0136-notes_left-L005}{}{}{\Indent (4) Cod. HDG.}
\end{minipage}
\hfill\begin{minipage}[t]{.48\textwidth}\vspace{0pt}\fontsize{9}{11.5}\selectfont\setlength{\GutterOffset}{0pt}
\hypertarget{AB01-PDF0136-N05}{}
\DLine{AB01-PDF0136-notes_right-L001}{}{}{\Indent (5) Cod. HS.}
\hypertarget{AB01-PDF0136-N06}{}
\DLine{AB01-PDF0136-notes_right-L002}{}{}{\Indent (6) Secundae apud Platonem $55''$.}
\hypertarget{AB01-PDF0136-N07}{}
\DLine{AB01-PDF0136-notes_right-L003}{}{}{\Indent (7) Plato $58^p15'34''$. Cfr. adnot. {\itshape e} huius capitis.}
\hypertarget{AB01-PDF0136-N08}{}
\DLine{AB01-PDF0136-notes_right-L004}{}{}{\Indent (8) Cod. $1^\circ4'24''$; Plato $1^\circ4'54''$.}
\hypertarget{AB01-PDF0136-N09}{}
\DLine{AB01-PDF0136-notes_right-L005}{}{}{\Indent (9) Cod. $1^\circ1'24''$; Plato $31^\circ4'54''$.}
\end{minipage}
